\input{preamble.tex}

\title{Отчет по лабораторной работе № 7. \\ Учёт физических параметров сети}
\author{Сущенко Алина Николаевна \\ НПИбд-01-23}
\date{2026}

\begin{document}

\maketitle
\newpage

\tableofcontents

\newpage
\section{Цель работы}
Приобретение практических навыков по работе с физической рабочей областью Packet Tracer, а также учёту физических параметров сети.

\newpage
\section{Выполнение работы}

\subsection{Настройка физической рабочей области}
\begin{enumerate}
\item В физической рабочей области Packet Tracer присвоили название городу - \texttt{ansusenko} (Рис. \ref{01}).
\begin{center}
    \centering
    \includegraphics[width=\textwidth]{../images/01.png}
    \captionof{figure}{Присвоение названия городу Moscow.}
    \label{01}
\end{center}

\item Щёлкнув на изображении города, перешли на уровень зданий. Присвоили зданию название \texttt{ansusenko}. Добавили здание для территории \texttt{ansusenko-p} (Рис. \ref{02}).
\begin{center}
    \centering
    \includegraphics[width=\textwidth]{../images/02.png}
    \captionof{figure}{Здания ansusenko и ansusenko-p.}
    \label{02}
\end{center}

\item Щёлкнув на изображении здания ansusenko, переместили изображение, обозначающее серверное помещение (Main Wiring Closet), в него (Рис. \ref{03}).
\begin{center}
    \centering
    \includegraphics[width=\textwidth]{../images/03.png}
    \captionof{figure}{Размещение серверного помещения в здании ansuenko.}
    \label{03}
\end{center}

\item Внутри серверного помещения разместили сетевое оборудование: коммутаторы \texttt{msk-ansusenko-sw-1}, \texttt{msk-ansusenko-sw-01}, \texttt{msk-ansusenko-sw-2}, \texttt{msk-ansusenko-sw-3}, \texttt{msk-ansusenko-sw-4}, а также серверы \texttt{ansusenko-web}, \texttt{ansusenko-file}, \texttt{ansusenko-mail} (Рис. \ref{04}).
\begin{center}
    \centering
    \includegraphics[width=\textwidth]{../images/04.png}
    \captionof{figure}{Размещение сетевого оборудования в серверной стойке.}
    \label{04}
\end{center}

\item На территории \texttt{ansuenko-p} разместили коммутатор \texttt{msk-ansusenko-sw-1} и оконечные устройства \texttt{other-ansusenko-1}, \texttt{dk-ansusenko-1} (Рис. \ref{05}).
\begin{center}
    \centering
    \includegraphics[width=\textwidth]{../images/05.png}
    \captionof{figure}{Размещение оборудования на территории ansusenko-p.}
    \label{05}
\end{center}
\end{enumerate}

\subsection{Настройка IP-адресации}
\begin{enumerate}
\item На коммутаторе \texttt{msk-ansusenko-sw-01} выполнили команду ping с IP-адресом \texttt{10.128.1.5} прежде, чем отдалили на расстояние 1000 м. (Рис. \ref{06}).
\begin{center}
    \centering
    \includegraphics[width=\textwidth]{../images/06.png}
    \captionof{figure}{Пинг на коммутаторе msk-ansusenko-sw-01.}
    \label{06}
\end{center}

\item На коммутаторе \texttt{msk-ansusenko-sw-1} выполнили команду ping с IP-адресом \texttt{10.128.1.6} после того, как отдалили на расстояние 1000 м.(Рис. \ref{07}).
\begin{center}
    \centering
    \includegraphics[width=\textwidth]{../images/07.png}
    \captionof{figure}{Неудачный пинг.}
    \label{07}
\end{center}
\end{enumerate}

\item Выполнили ping с коммутатора \texttt{msk-ansusenko-sw-01} на коммутатор \texttt{msk-ansusenko-sw-1} по IP-адресу \texttt{10.128.1.6} (Рис. \ref{08}). Убедились, что пинг недоступен.:
\end{enumerate}

\begin{center}
    \centering
    \includegraphics[width=\textwidth]{../images/08.png}
    \captionof{figure}{Неудачный пинг.}
    \label{08}
\end{center}

\subsection{Включение учёта физических характеристик среды передачи}
\begin{enumerate}
\item В меню \texttt{Options} $\rightarrow$ \texttt{Preferences} во вкладке \texttt{Interface} активировали разрешение на учёт физических характеристик среды передачи (\texttt{Enable Cable Length Effects}).
\end{enumerate}

\item В рабочее пространство добавили повторители (Repeater-PT). Присвоили им соответствующие названия msk-ansuseko-mc-1 и msk-ansusenko-mc-01. Заменим имеющиеся модули на PT-REPEATER-NM-1FFE и PT-REPEATER-NM-1CFE для подключения оптоволокна и витой пары по технологии Fast Ethernet(Рис. \ref{09}):
\end{enumerate}

\begin{center}
    \centering
    \includegraphics[width=\textwidth]{../images/09.png}
    \captionof{figure}{Вид повторителя}
    \label{09}
\end{center}

\item Подключили повторители по схеме(Рис.\ref{10}):
\end{enumerate}

\begin{center}
    \centering
    \includegraphics[width=\textwidth]{../images/10.png}
    \captionof{figure}{Новый вид схемы (отрезок с повторителями)}
    \label{10}
\end{center}

\subsection{Изменение расстояния между территориями}
\begin{enumerate}
\item В физической рабочей области Packet Tracer разместили две территории на расстоянии более 100 м друг от друга. Фактическое расстояние составило \textbf{991.37 метров} (Рис. \ref{11}).
\begin{center}
    \centering
    \includegraphics[width=\textwidth]{../images/11.png}
    \captionof{figure}{Расстояние между территориями 991.37 м, соеденины.}
    \label{11}
\end{center}

\begin{enumerate}
\item Проверили доступ между городами.(Рис. \ref{12}, \ref{13}).
\begin{center}
    \centering
    \includegraphics[width=\textwidth]{../images/12.png}
    \captionof{figure}{Ping от ansusenko к ansusenko-p.}
    \label{12}
\end{center}

\begin{center}
    \centering
    \includegraphics[width=\textwidth]{../images/13.png}
    \captionof{figure}{Ping от ansusenko-p к ansusenko.}
    \label{13}
\end{center}

\section{Выводы}
В результате выполнения лабораторной работы приобрели практические навыки по работе с физической рабочей областью Packet Tracer.
\end{document}
